% Additions to the preliminary results report (2026-10-29).
% 1. Overview table with the added metrics. 2. Its metric descriptions.
% 3. Table of the major changes since the first version, with explanations.
% Requires \usepackage{booktabs,multirow,amsmath,array,float}.

\begin{table}[H]
\centering
\small
\setlength{\tabcolsep}{4pt}
\caption{Overview of the four fronts. Comparison 1 (columns 3--4): standard backprop vs our dense configuration (no energy penalty). Comparison 2 (columns 5--6): backprop trained with the same energy term our method pays ($\mu = 10^{-4}$; its Adam learning rate tuned over $\{3\cdot10^{-4}, 10^{-3}, 3\cdot10^{-3}\}$, best shown) vs our sparse configuration at the same $\mu$. MNIST, 784--128--10, batch 256, mean $\pm$ std over 3 seeds. Training step, step energy and step memory from the controlled benchmark (Tables~\ref{tab:t1-bench} and \ref{tab:t2-bench}); training time per epoch, total training energy and time to 95\% measured during the training runs. Bold: better mean within each comparison. $^\dagger$Difference smaller than one standard deviation in at least one comparison.}
\label{tab:overview}
\begin{tabular}{>{\raggedright\arraybackslash}p{2.1cm}>{\raggedright\arraybackslash}p{3.5cm}cccc}
\toprule
Category & Metric & Backprop & Ours (dense) & Backprop + energy & Ours (sparse) \\
\midrule
\multirow{2}{*}{\parbox[c]{2.1cm}{\raggedright\textit{Accuracy}}} & Test accuracy (\%)$^\dagger$ & $\boldsymbol{97.63 \pm 0.05}$ & $97.59 \pm 0.09$ & $\boldsymbol{97.12 \pm 0.07}$ & $97.06 \pm 0.15$ \\
 & ECE (calibration) & $\boldsymbol{0.0034 \pm 0.0002}$ & $0.0038 \pm 0.0001$ & $\boldsymbol{0.0044 \pm 0.0010}$ & $0.0087 \pm 0.0007$ \\
\midrule
\multirow{5}{*}[-30pt]{\parbox[c]{2.1cm}{\raggedright\textit{Energy}}} & Synaptic operations per image, network$^\dagger$ & $11339 \pm 243$ & $\boldsymbol{3858 \pm 474}$ & $525 \pm 120$ & $\boldsymbol{472 \pm 15}$ \\
 & Synaptic operations per image, total$^\dagger$ & $343802 \pm 255$ & $\boldsymbol{336322 \pm 466}$ & $332989 \pm 128$ & $\boldsymbol{332936 \pm 24}$ \\
 & GPU energy per training step (mJ) & $991 \pm 19$ & $\boldsymbol{768 \pm 11}$ & $1078 \pm 4$ & $\boldsymbol{682 \pm 3}$ \\
 & Total training GPU energy (kJ) & $17.78 \pm 1.23$ & $\boldsymbol{14.10 \pm 0.15}$ & $23.10 \pm 7.57$ & $\boldsymbol{14.21 \pm 1.43}$ \\
 & Inference GPU energy per image (mJ)$^\dagger$ & $\boldsymbol{1.48 \pm 0.13}$ & $1.49 \pm 0.03$ & $\boldsymbol{1.42 \pm 0.03}$ & $1.44 \pm 0.06$ \\
\midrule
\multirow{5}{*}[-24pt]{\parbox[c]{2.1cm}{\raggedright\textit{Efficiency}}} & Training step (ms) & $7.64 \pm 0.05$ & $\boldsymbol{5.53 \pm 0.16}$ & $9.21 \pm 0.01$ & $\boldsymbol{5.37 \pm 0.03}$ \\
 & Training time per epoch (s) & $4.57 \pm 0.19$ & $\boldsymbol{3.82 \pm 0.24}$ & $5.74 \pm 1.43$ & $\boldsymbol{3.73 \pm 0.13}$ \\
 & Epochs to 95\% validation accuracy$^\dagger$ & $2.0 \pm 1.0$ & $\boldsymbol{1.7 \pm 0.6}$ & $3.0 \pm 0.0$ & $\boldsymbol{1.3 \pm 0.6}$ \\
 & Training time to 95\% validation accuracy (s)$^\dagger$ & $9.4 \pm 4.3$ & $\boldsymbol{6.7 \pm 1.9}$ & $17.8 \pm 5.0$ & $\boldsymbol{5.8 \pm 1.8}$ \\
 & Peak GPU memory, training step (MiB) & $\boldsymbol{230.9 \pm 0.0}$ & $271.6 \pm 0.0$ & $\boldsymbol{234.1 \pm 0.0}$ & $248.5 \pm 0.0$ \\
\midrule
\multirow{3}{*}[-18pt]{\parbox[c]{2.1cm}{\raggedright\textit{Explainability and robustness}}} & Faithfulness margin, 25\% removed & $0.04 \pm 0.02$ & $\boldsymbol{0.40 \pm 0.15}$ & $\boldsymbol{0.41 \pm 0.07}$ & $0.10 \pm 0.02$ \\
 & Credit vs. single-neuron removal (Spearman) & $0.57 \pm 0.03$ & $\boldsymbol{0.81 \pm 0.11}$ & $\boldsymbol{0.93 \pm 0.03}$ & $0.84 \pm 0.04$ \\
 & Accuracy after removing a random 25\% (\%)$^\dagger$ & $\boldsymbol{92.8 \pm 1.0}$ & $81.3 \pm 6.1$ & $89.5 \pm 2.2$ & $\boldsymbol{91.4 \pm 0.4}$ \\
\midrule
\textit{Backprop. Check} & Autograd calls in training (per seed) & 10{,}700 & 0 & 10{,}700 & 0 \\
\bottomrule
\end{tabular}
\end{table}

{\small
\noindent\textbf{Overview of the four fronts}\par\smallskip
\noindent\textbf{Accuracy:}
\begin{itemize}
  \item \textbf{Test accuracy (\%):} The share of the 10,000 test images classified correctly, using the checkpoint with the best validation accuracy. The prediction is the read-out neuron with the most accumulated membrane potential over the 25 time steps. Higher is better.
  \item \textbf{ECE (calibration):} Expected calibration error. Test predictions are grouped into 15 bins by confidence; in each bin, the gap between average confidence and actual accuracy is taken, and the gaps are averaged, weighted by bin size. It measures whether a model's confidence matches how often it is right (0 = perfect). Probabilities are temperature-scaled on the validation set first, so both models are compared at their best calibration. Lower is better.
\end{itemize}

\noindent\textbf{Energy:}
\begin{itemize}
  \item \textbf{Synaptic operations per image, network:} How many synapses are activated by hidden-layer spikes while classifying one test image: a spike of a neuron with $n$ outgoing synapses counts $n$. It is the work a neuromorphic chip would do, which only computes when a neuron fires. Lower is better.
  \item \textbf{Synaptic operations per image, total:} The same, including the input layer's spikes. The input is identical for both models and dominates this count in a network with one hidden layer, so the difference here is small even when the network count differs severalfold. Lower is better.
  \item \textbf{GPU energy per training step (mJ):} Energy the GPU draws for one training step (forward pass plus the model's own update; for backprop this includes its backward pass), read from the GPU's energy counter in the controlled benchmark. Lower is better.
  \item \textbf{Total training GPU energy (kJ):} GPU energy over the whole 50-epoch training run, including validation passes. It reflects the total compute spent to train a model. Lower is better.
  \item \textbf{Inference GPU energy per image (mJ):} GPU energy to classify one test image, from the controlled benchmark. On a GPU, which computes densely, sparsity does not lower this; the synaptic-operation counts show the hardware-independent difference. Lower is better.
\end{itemize}

\noindent\textbf{Efficiency:}
\begin{itemize}
  \item \textbf{Training step (ms):} Time for one training step on a batch of 256 images, from the controlled benchmark (same preloaded data, models interleaved, repeated). Lower is better.
  \item \textbf{Training time per epoch (s):} Time for one pass over the 55,000 training images during the actual training runs, excluding validation. Lower is better.
  \item \textbf{Epochs to 95\% validation accuracy:} After every epoch, the model is evaluated on the 5,000 validation images. This is the first epoch at which that accuracy reaches 95\%, averaged over seeds. It measures how few passes over the training data a model needs to become useful. Lower is better.
  \item \textbf{Training time to 95\% validation accuracy (s):} The total time spent training (forward passes and weight updates on the training images) from the start up to the end of that epoch. The validation evaluations used to check the accuracy are not counted, so the number is the cost of learning only. It combines how many epochs are needed with how long each one takes. Lower is better.
  \item \textbf{Peak GPU memory, training step (MiB):} The most GPU memory in use during a training step in the controlled benchmark, including the preloaded batches. Our method keeps information about every time step of the window (for its local updates), which costs more memory than backprop's stored activations. Lower is better.
\end{itemize}

\noindent\textbf{Explainability and robustness:}
\begin{itemize}
  \item \textbf{Faithfulness margin, 25\% removed:} Each model assigns a credit score to each of its 128 hidden neurons. The margin is test accuracy after removing 32 random neurons minus test accuracy after removing the 32 top-credited ones (best of the three credit scores). Positive means the credit identifies the neurons that matter. It also rewards networks that concentrate their computation in a few neurons. Higher is better.
  \item \textbf{Credit vs.\ single-neuron removal (Spearman):} Each hidden neuron is removed on its own and the resulting increase in test loss is measured. The score is the rank correlation between those measured effects and the credit scores (1 = the credit orders the neurons exactly by their real importance). Unlike the margin, it does not depend on how concentrated the network is. Higher is better.
  \item \textbf{Accuracy after removing a random 25\% (\%):} Test accuracy after silencing 32 randomly chosen hidden neurons, averaged over five random draws. It measures robustness to losing neurons (for example, hardware faults). Higher is better.
\end{itemize}

\noindent\textbf{Backprop. check:}
\begin{itemize}
  \item \textbf{Autograd calls in training (per seed):} How many times PyTorch's backpropagation engine was called during training. Backprop calls it once per batch (214 batches $\times$ 50 epochs = 10,700). Ours must be 0; this verifies that our method never uses backprop.
\end{itemize}
}

\begin{table}[H]
\centering
\small
\setlength{\tabcolsep}{4pt}
\caption{Major changes since the first version of this report, in the order they were made, and their effect. Test accuracy on MNIST, 784--128--10. Each change is measured against the configuration just before it, on the same code and settings otherwise; ``3 seeds'' marks means over seeds 1--3, other rows are single runs (seed-to-seed noise $\approx 0.3$ points). The training-step times are measured under controlled conditions on the same GPU. Batch size 128 unless noted.}
\label{tab:changes}
\begin{tabular}{>{\raggedright\arraybackslash}p{4.3cm}ccc>{\raggedright\arraybackslash}p{4.4cm}}
\toprule
Change & Test acc.\ before (\%) & Test acc.\ after (\%) & $\Delta$ (points) & Other major effects \\
\midrule
First version (all proposal terms, fixed random feedback, spiking read-out) & -- & $94.19 \pm 0.08$ (3 seeds) & -- & Gap to backprop $-3.44$ points \\
\midrule
1. Leaky read-out with Kolen-Pollack feedback & $94.25$ & $95.62$ & $+1.37$ & Network synaptic operations $4{,}028 \to 986$ (4$\times$ fewer); faithfulness margin $+0.06 \to +0.20$ \\
2. Diversity term removed & $95.62$ & $96.21$ & $+0.59$ & Synaptic operations $986 \to 634$; faithfulness margin $+0.20 \to +0.61$; training step 11\% faster \\
3. Energy weight $\mu$ lowered to 0 & $96.20 \pm 0.03$ (3 seeds) & $97.56 \pm 0.07$ (3 seeds) & $+1.36$ & Synaptic operations $630 \to 2{,}470$ (still 4$\times$ fewer than backprop) \\
4. Batch size 256 (both models) & $97.58$ & $97.59$ & $+0.01$ & Training time per epoch $7.30 \to 4.67$ s ($-36\%$); the same for backprop \\
5. Faster implementation of the same update & -- & -- & $0$ (identical updates) & Training step $10.5 \to 5.9$ ms ($-44\%$); now faster than backprop (8.4 ms) \\
\midrule
6. Sparse configuration: lower threshold target with the energy term on & $95.70 \pm 0.15$ (3 seeds, $\mu = 10^{-4}$) & $97.09 \pm 0.07$ (3 seeds) & $+1.39$ & Same sparsity (467 synaptic operations); makes comparison 2 a tie instead of a loss \\
\midrule
Now (dense configuration, batch 256) & -- & $97.59 \pm 0.09$ (3 seeds) & $+3.40$ vs.\ first version & Gap to backprop $-0.03 \pm 0.07$ points; training energy $36.3 \to 14.1$ kJ \\
\bottomrule
\end{tabular}
\end{table}

{\small
\noindent\textbf{What each change is:}
\begin{itemize}
  \item \textbf{First version.} The full proposal: three-factor updates, diversity term $V$, metabolic and barrier energy $E$, Eligibility-Shapley credit modulating the learning signal, a spiking read-out layer, fixed random feedback matrices carrying the output error to the hidden layer, and plain SGD (learning rate 2.0). Before change 1, two smaller steps had no effect on accuracy: switching to SGD with momentum 0.9 (learning rate 0.2) made training converge sooner, and the credit modulation was turned off because it changed nothing (94.20\% with and without).
  \item \textbf{1. Leaky read-out with Kolen-Pollack feedback.} The output neurons no longer spike: each one integrates its input as a leaky membrane potential, and the prediction is the largest accumulated potential. This makes the output layer's own update exact. The hidden layer's learning signal travels through a feedback matrix that receives the same local update as the output weights and decays with them (Kolen-Pollack), so it converges to the output weights within about two epochs instead of staying random. No backward pass is involved; each matrix changes only through its own local update. Either change alone hurt; together they gave the largest single gain. The energy weight was reduced 4$\times$ at the same time, because the leaky read-out's task loss is smaller and the old weight silenced the hidden layer.
  \item \textbf{2. Diversity term removed.} With the leaky read-out and Kolen-Pollack feedback, the diversity term (which rewards neurons for firing differently from each other) no longer helped. Removing it improved accuracy, sparsity, credit faithfulness and speed at once.
  \item \textbf{3. Energy weight $\mu$ lowered to 0.} The metabolic energy term charges every spike. A sweep over 3 seeds showed it trades accuracy for sparsity monotonically rather than improving accuracy. Setting it to 0 gives the most accurate configuration, which is still 3--4$\times$ sparser than backprop. A firing-rate floor and adaptive thresholds were added in the same batch; they keep neurons from going silent and matter more on Fashion-MNIST, but did not change MNIST accuracy.
  \item \textbf{4. Batch size 256.} Training on 256 images per step instead of 128, for both models. Accuracy is unchanged for both and each epoch takes about a third less time.
  \item \textbf{5. Faster implementation.} The same update is now computed with far fewer GPU operations: one quantity (the surrogate derivative) is computed once per time window instead of at every time step, and all logged values are copied from the GPU in one transfer instead of about eleven. Both models were limited by the overhead of launching many small GPU operations, not by the arithmetic. The resulting weights are identical.
  \item \textbf{6. Lower threshold target with the energy term on.} Each hidden neuron adapts its firing threshold toward a target firing rate (5\%). When the energy term is on, it pushes firing rates below that target, so the two mechanisms work against each other. Turning the adaptive thresholds off (or lowering the target to 2\%) removes the conflict. This only concerns the sparse configuration of comparison 2.
\end{itemize}
}
